\documentclass{article}
\PassOptionsToPackage{round,authoryear}{natbib}
\usepackage[dblblindworkshop]{neurips_2026}
\workshoptitle{3rd Workshop on Attributing Model Behavior at Scale: Data Attribution and Provenance (ATTRIB 2026)}

\usepackage[utf8]{inputenc}
% NOTE: no [T1]{fontenc}. With T1 the EC typewriter falls back to a bitmap PK font, which
% lands in the PDF as Type 3; the NeurIPS instructions require Type 1 or embedded TrueType
% only. OT1 keeps the template's Times body font and gives a real Type 1 CMTT for \texttt.
\usepackage{hyperref}
\usepackage{url}
\usepackage{booktabs}
\usepackage{amsmath,amssymb}
\usepackage{amsfonts}
\usepackage{amsthm}
\usepackage{microtype}
\usepackage{graphicx}

\newtheorem{lemma}{Lemma}

\title{Model Family Is Recoverable From Co-Failure Alone: Population-Scale Provenance Signal Without Training-Data Access}

\author{
  Anonymous Author(s) \\
  Affiliation \\
  Address \\
  \texttt{email}
}

\begin{document}

\maketitle

\begin{abstract}
Contributive and corroborative attribution answer ``which training data explains this output,''
one output at a time, and need access to training data or model internals that is unavailable for
most open leaderboard submissions. We ask whether anything about provenance survives in the only
signal such a population does expose: which items each model gets wrong. Conditioning a
model-by-item failure matrix on the exact, fit-free null implied by Rasch sufficiency and taking
the leading eigenvectors of the residual correlation matrix, a $5$-nearest-neighbour classifier
recovers a model's base family with leave-one-out accuracy $0.807$ over $916$ labelled models from
five families, against a label-permutation null of $0.287 \pm 0.017$ and a majority-class rate of
$0.382$. The same pipeline run on a draw from the exact null returns $0.318$, at chance, so the
construction does not manufacture the signal. It is not an accuracy artifact: projecting model
accuracy out of the loadings leaves accuracy at $0.809$. It is partly, but not mostly, a
near-duplicate artifact: deduplicating at pairwise agreement $0.95$ leaves $0.787$, and the much
more aggressive $0.90$ threshold leaves $0.551$, still far above chance. We report one negative
result prominently: margin conditioning is not what creates the signal, since unconditioned
correlation eigenvectors already reach $0.760$. We did not set out to build an attribution method,
and this is not one; it is evidence that population-scale behavioural correlation carries
recoverable provenance information, obtained on a population where a pre-registered lineage
analysis had to be abandoned for lack of declared metadata.
\end{abstract}

\section{The gap we are pointing at}

Contributive attribution estimates the causal effect of specific training examples on a specific
output \citep{kohliang2017,park2023trak}; corroborative attribution finds sources that support a
specific output \citep{worledge2024}. Both operate on one model and one output, and both need
access, to training data or to internals, that an open leaderboard population does not provide.
Thousands of models are submitted with undisclosed lineage. The question there is not ``which
example explains this generation'' but ``which models share unstated ancestry, and can that be seen
at all from behaviour.''

We met this question from the measurement side, while characterising correlated failure between
open language models, and we met it as a dead end: a pre-registered lineage decomposition had to be
abandoned because the metadata it needed does not exist at usable coverage
(Section~\ref{sec:opacity}). This paper reports what we did instead, and what it found.

\section{Setup: the exact conditional null and its residual spectrum}

Let $F\in\{0,1\}^{N\times M}$ with $F_{im}=1$ iff model $i$ fails item $m$, row sums $r_i$, column
sums $c_m$.

\begin{lemma}[Mean co-failure depends only on item margins]
\label{lem:mean}
$O=\frac{1}{MN(N-1)}\sum_m(c_m^2-c_m)$, so any randomization preserving the item margins leaves the
mean pairwise co-failure rate exactly invariant, with zero variance.
\end{lemma}
\begin{proof}
$\sum_{i\neq j}F_{im}F_{jm}=(\sum_i F_{im})^2-\sum_i F_{im}^2=c_m^2-c_m$ for $F_{im}\in\{0,1\}$.
\end{proof}

Lemma~\ref{lem:mean} says the first thing one would compute carries no provenance information at
all once item difficulty is accounted for; the same degeneracy is known in ecology and
psychometrics \citep{schluter1984,gotelli2000,cronbach1951}. We therefore work with second-order
structure. The row and column margins of $F$ are the jointly sufficient statistics of the Rasch
model, so the conditional law of $F$ given its margins is uniform on the margin-matched set for
every parameter value \citep{rasch1960,andersen1973}; we sample it with curveball
\citep{strona2014,carstens2015}. Writing $P$ for the Rasch mean field matching the margins,
$D=(FF^{\!\top}-PP^{\!\top})/M$ and $R$ for $D$ scaled to unit diagonal, we take the top-$k$
eigenvectors of $R$, with $k$ the number of eigenvalues exceeding the exact null's spectral edge.

\paragraph{Population.} A dated snapshot of the archived Open LLM Leaderboard, ARC-Challenge:
$N=3{,}762$ models, $M=1{,}169$ items, read from result files at zero inference cost. On this
population $k=4$ (null edge $79.5$) and the residual correlation is $5.4\times$ its null. This is a
larger, later harvest than the $N=1{,}362$ sample in our companion manuscript, and every number
here refers to $N=3{,}762$. The archive grows continuously and a scheduled process keeps harvesting
it, so the snapshot is pinned by content hash (SHA-256 \texttt{79d2490e\ldots}) in the released
artifacts rather than described as current; re-running on a later snapshot will not reproduce these
counts exactly, and the pin is what makes the run reproducible.

\paragraph{Labels.} Base family is assigned by string matching on the model identifier against a
list of known base-model names. This is a coarse lower bound on relatedness and explicitly not
lineage ground truth; it attributes $34.5$--$35.7\%$ of identifiers per benchmark. We use the five
largest families, $916$ labelled models: Mistral $350$, Llama-2 $211$, Llama-3 $188$, Qwen $84$,
Mixtral $83$.

\section{Result: family is recoverable from co-failure}
\label{sec:result}

Figure~\ref{fig:family} shows the effect directly. Leave-one-out $5$-nearest-neighbour
classification in the standardised $k$-dimensional loading space recovers family at accuracy $\mathbf{0.807}$, against a label-permutation null of
$0.287\pm0.017$ (1{,}000 permutations, $p=0.001$, the resolution floor) and a majority-class rate
of $0.382$. Recall is above chance for four of five families: Mistral $0.926$, Llama-2 $0.839$,
Llama-3 $0.777$, Qwen $0.714$, Mixtral $0.386$. Accuracy rises monotonically with the number of
eigenvectors used, from $0.452$ at $k=1$ to $0.864$ at $k=6$ and $0.905$ at $k=20$, so the signal
is distributed across many residual dimensions rather than carried by one.

Every kill condition and control was specified before the experiment ran.

\begin{figure}[t]
\centering
\includegraphics[width=\linewidth]{../../../figures/fig_w1_family.pdf}
\caption{Models projected on the top two eigenvectors of the margin-conditioned residual
correlation matrix, coloured by base family (left), and the identical pipeline applied to a
curveball replicate of the same matrix (right). The replicate preserves both margins exactly and so
destroys any real association between behaviour and family: it is what the left panel would look
like if the construction were manufacturing the structure. Only the five largest families are
drawn; the classifier of Section~\ref{sec:result} uses four eigenvectors, not the two shown here.}
\label{fig:family}
\end{figure}

\begin{table}[t]
\centering
\small
\begin{tabular}{llrrl}
\toprule
Arm & What it tests & Acc. & Perm. & Reading \\
\midrule
Main & residual eigenvectors, real labels & \textbf{0.807} & $0.287$ & signal present \\
Exact-null draw & pipeline on a curveball replicate & $0.318$ & $0.289$ & at chance \\
Accuracy only & accuracy as the sole feature & $0.432$ & $0.288$ & carries some, not most \\
Accuracy proj.\ out & loadings orthogonalised to accuracy & $0.809$ & $0.286$ & not an accuracy artifact \\
Raw correlation & unconditioned eigenvectors & $0.760$ & $0.287$ & conditioning not the source \\
Dedup $0.95$ & $N=2{,}350$, $591$ labelled & $0.787$ & $0.276$ & survives \\
Dedup $0.90$ & $N=1{,}288$, $301$ labelled & $0.551$ & $0.264$ & weakened, above chance \\
\bottomrule
\end{tabular}
\caption{Pre-registered controls. The exact-null arm (KW3) is the calibration check: because the
curveball replicate destroys any real association between behaviour and family while preserving
both margins, a pipeline that reported signal there would be reporting an artifact. It does not.}
\label{tab:controls}
\end{table}

\paragraph{What we are careful not to claim.} Three things in Table~\ref{tab:controls} cut against
the strongest reading, and we state them rather than bury them. First, margin conditioning is
\emph{not} what produces the signal: unconditioned correlation eigenvectors already reach $0.760$,
so conditioning adds about five points, not the effect. The exact null earns its place here as the
calibration arm and as what makes the mean-level statistic provably uninformative
(Lemma~\ref{lem:mean}), not as the source of the result. Second, part of the signal is
near-duplicate structure: accuracy falls from $0.807$ to $0.551$ under aggressive deduplication at
agreement $0.90$, which discards two thirds of the population. The residue is still roughly twice
the permutation null, so the signal is not only duplicates, but it is partly duplicates. Third, the
labels are name strings. A model named for a base it was not derived from, or derived from a base it
is not named for, is mislabelled here, and we have no way to audit that.

\section{The provenance opacity this was a response to}
\label{sec:opacity}

We did not choose an indirect route for its own sake. The direct one failed, and quantifiably. The
leaderboard's optional, self-reported \texttt{base\_model} field resolved a declared parent for
$317$ of $1{,}362$ models ($23.3\%$) in our earlier harvest, with $315$ repositories no longer
retrievable at all. A $40\%$ coverage threshold had been registered in advance as the point below
which a lineage decomposition would not be trusted; it fired, and the decomposition was dropped
rather than run on the surviving quarter. Name-string matching, used above, reaches $34.5$--$35.7\%$
and is coarser still.

That failure is the reason this result matters for this workshop rather than only for the
correlated-failure literature. On a population where declared provenance is missing for three
quarters of models, a behavioural signal that recovers family at $0.81$ from outcome data alone is
the only provenance evidence available. It costs one Rasch fit and an eigendecomposition, needs no
access to weights, training data, or inference, and scales to the full population.

\section{What this is and is not}

It is not an attribution method. It returns a coordinate in a residual eigenspace, not a claim
about which examples caused which behaviour, and it has no notion of a specific output to explain.
The natural next questions are whether the loadings support pairwise ancestry tests rather than
family classification, whether the recovered structure survives against config-hash ground truth
instead of name strings, and whether it degrades gracefully as families proliferate. What the
result does establish is the premise those questions need: population-scale co-failure structure is
not provenance-blind. Under an exact null that removes item difficulty by construction, and after
projecting out model accuracy, behaviour alone still locates a model in its family most of the time.

\section{Limitations}

Single benchmark, single archive, one coarse labelling scheme; the five families used are the five
largest and the sixth-largest onward are excluded. Deduplication at $0.90$ costs the signal a
third of its margin over chance, so a population with a different duplicate structure could behave
differently. We report classification accuracy, not a calibrated provenance probability, and offer
no decision procedure. An alternative explanation we cannot exclude is that the residual dimensions
track training recipe or data mixture, which correlates with family without ancestry: that would
make this a signature of convergent practice rather than of lineage, and distinguishing the two
needs the ground truth this population lacks. Code, data, the pre-registration with its dated kill
conditions, and every artifact behind the numbers above are archived and will be linked in the
camera-ready. A large language model was used as a coding and drafting assistant; the authors
specified and pre-registered the hypotheses and kill conditions, verified every derivation and
every number against executed-run artifacts, and are responsible for the content.

\begin{thebibliography}{99}
\small
\bibitem[Andersen(1973)]{andersen1973} E.~B. Andersen. A goodness of fit test for the Rasch model. \emph{Psychometrika}, 38:123--140, 1973.
\bibitem[Carstens(2015)]{carstens2015} C.~J. Carstens. Proof of uniform sampling of binary matrices with fixed row sums and column sums for the fast Curveball algorithm. \emph{Physical Review E}, 91:042812, 2015.
\bibitem[Cronbach(1951)]{cronbach1951} L.~J. Cronbach. Coefficient alpha and the internal structure of tests. \emph{Psychometrika}, 16:297--334, 1951.
\bibitem[Gotelli(2000)]{gotelli2000} N.~J. Gotelli. Null model analysis of species co-occurrence patterns. \emph{Ecology}, 81(9):2606--2621, 2000.
\bibitem[Koh \& Liang(2017)]{kohliang2017} P.~W. Koh and P.~Liang. Understanding black-box predictions via influence functions. \emph{ICML}, 2017.
\bibitem[Park et al.(2023)]{park2023trak} S.~M. Park, K.~Georgiev, A.~Ilyas, G.~Leclerc, and A.~Madry. TRAK: Attributing model behavior at scale. \emph{ICML}, 2023.
\bibitem[Rasch(1960)]{rasch1960} G.~Rasch. \emph{Probabilistic Models for Some Intelligence and Attainment Tests}. Danmarks P\ae dagogiske Institut, 1960.
\bibitem[Schluter(1984)]{schluter1984} D.~Schluter. A variance test for detecting species associations, with some example applications. \emph{Ecology}, 65:998--1005, 1984.
\bibitem[Strona et al.(2014)]{strona2014} G.~Strona, D.~Nappo, F.~Boccacci, S.~Fattorini, and J.~San-Miguel-Ayanz. A fast and unbiased procedure to randomize ecological binary matrices with fixed row and column totals. \emph{Nature Communications}, 5:4114, 2014.
\bibitem[Worledge et al.(2024)]{worledge2024} T.~Worledge, J.~H. Shen, N.~Meister, C.~Winston, and C.~Guestrin. Unifying corroborative and contributive attributions in large language models. \emph{IEEE Conference on Secure and Trustworthy Machine Learning (SaTML)}, 2024.
\end{thebibliography}

\end{document}
